% SEGMENT OLP-0007-S01
% Part: sets-functions-relations
% Chapter: sets
% Section: important-sets

\documentclass[../../../include/open-logic-section]{subfiles}

\begin{document}

\olfileid{sfr}{set}{imp}
\olsection{முக்கியமான சில கணங்கள்}

\begin{ex}
நாம் பெரும்பாலும் கணிதப் பொருள்களை !!{element}s ஆகக் கொண்ட கணங்களையே
கையாளுவோம். அவற்றில் நான்கு கணங்கள் தனிப்பெயர்கள் பெறும் அளவுக்கு முக்கியமானவை:
\begin{multline*}
    \Nat = \{0, 1, 2, 3, \ldots\} \\
    \shoveright{\text{இயல் எண்களின் கணம்}}\\
    \shoveleft{\Int = \{\ldots, -2, -1, 0, 1, 2, \ldots\}} \\
    \shoveright{\text{முழுக்களின் கணம்}}\\
    \shoveleft{\Rat = \Setabs{\nicefrac{m}{n}}{m, n \in \Int\text{ மற்றும் }n \neq 0}}\\
    \shoveright{\text{விகிதமுறு எண்களின் கணம்}}\\
    \shoveleft{\Real = (-\infty, \infty)}\\
    \text{மெய்யெண்களின் கணம் (தொடர்ச்சித் தொகுதி)}
\end{multline*}
இவை அனைத்தும் \emph{முடிவுறாக்} கணங்கள்; அதாவது, இவை ஒவ்வொன்றிலும்
முடிவிலா எண்ணிக்கையில் !!{element}s உள்ளன.

% SEGMENT OLP-0007-S02
இந்தக் கணங்களின் வழியே செல்லச்செல்ல, நம்மிடம் உள்ள எண்களுடன்
\emph{மேலும்} எண்களைச் சேர்க்கிறோம். உண்மையில்,
$\Nat \subseteq \Int \subseteq \Rat \subseteq \Real$ என்பது தெளிவாக இருக்க வேண்டும்:
ஒவ்வோர் இயல் எண்ணும் ஒரு முழு; ஒவ்வொரு முழுவும் ஒரு விகிதமுறு எண்;
ஒவ்வொரு விகிதமுறு எண்ணும் ஒரு மெய்யெண்.
அதேபோல், $\Nat \subsetneq \Int \subsetneq \Rat$ என்பதும் தெளிவாக இருக்க வேண்டும்;
ஏனெனில் $-1$ ஒரு முழு, ஆனால் இயல் எண் அல்ல;
$\nicefrac{1}{2}$ ஒரு விகிதமுறு எண், ஆனால் முழு அல்ல.
$\Rat \subsetneq \Real$ என்பது அவ்வளவு வெளிப்படையானதல்ல;
அதாவது, விகிதமுறு எண்களாக இல்லாத மெய்யெண்கள் சில உள்ளன
\oliflabeldef{sfr:arith:real:realline}{; இதனை \olref[arith][real]{realline} இல் மீண்டும் காண்போம்}{}.

% SEGMENT OLP-0007-S03
சில சமயங்களில் நேர்ம முழுக்களின் கணமான
$\PosInt = \{1, 2, 3, \dots\}$ என்பதையும், முதல் இரண்டு இயல் எண்களை மட்டும்
கொண்ட கணமான $\Bin = \{0, 1\}$ என்பதையும் பயன்படுத்துவோம்.
\end{ex}

% SEGMENT OLP-0007-S04
\begin{tagblock}{compsci}
\begin{ex}[சரங்கள்]
மற்றொரு சுவையான எடுத்துக்காட்டு, $A$ என்னும் எழுத்துக்கணத்தின் மீதான
\emph{முடிவுறு சரங்கள்} அனைத்தையும் கொண்ட $A^{*}$ என்னும் கணமாகும்:
$A$ இன் உறுப்புகளால் ஆன எந்த ஒரு முடிவுறு தொடர்முறையும் $A$ இன் மீதான ஒரு சரமாகும்.
ஒவ்வொரு எழுத்துக்கணம் $A$ க்கும், $A$ இன் மீதான சரங்களில்
\emph{வெற்றுச் சரம் $\Lambda$} என்பதையும் சேர்த்துக்கொள்கிறோம்.
எடுத்துக்காட்டாக,
\begin{multline*}
\Bin^*
=\{\Lambda,0,1,00,01,10,11,\\
000,001,010,011,100,101,110,111,0000,\ldots\}.
\end{multline*}
$x=x_{1}\ldots x_{n}\in A^{*}$ என்பது $A$ இன் $n$ ``எழுத்துகளால்'' ஆன
ஒரு சரம் எனில், அந்தச் சரத்தின் \emph{நீளம்} $n$ எனக் கூறி,
$\len{x}=n$ என எழுதுகிறோம்.
\end{ex}
\end{tagblock}

% SEGMENT OLP-0007-S05
\begin{ex}[முடிவுறாத் தொடர்முறைகள்]
எந்த ஒரு கணம் $A$ க்கும், $A$ இன் !!{element}s கொண்டு அமையும் முடிவுறாத்
தொடர்முறைகளின் கணமான $A^\omega$ என்பதையும் கருதலாம்.
$a_1a_2a_3a_4\dots$ என்னும் முடிவுறாத் தொடர்முறை என்பது
ஒரு திசையில் முடிவின்றி நீளும் பொருள்களின் பட்டியலாகும்;
அவற்றில் ஒவ்வொன்றும் $A$ இன் !!a{element} ஆகும்.
\end{ex}

\end{document}
